\section{Methods}\label{methods}

The dataset will be derived from the Caravan v1.6 CSV release,
specifically looking at US basins (camels), Great Britain basins
(camelsgb), and Australia basins (camelsaus), totaling 1760 basins.
Measurments for all basins will be daily from 1981-2014. Basins will be
kept for analysis based on overall quality and completeness.

Static attributes for the kept basins will represent the physical and
environmental characteristics of each basin. These attributes will be
standardized prior to clustering and model development, with all
preprocessing parameters estimated exclusively from the training
catchments within each validation fold.

To prevent information leakage, catchment attributes derived directly
from observed streamflow or other discharge-related measurements will be
excluded from the predictor set. The remaining static catchment
attributes will be examined for missing values, inconsistencies, and
potential redundancy before model development.

Meteorological predictors will be constructed using fixed time intervals
of 7, 30, 90, and 365 days. These intervals will capture short-term and
longer-term hydrological and climatic conditions that may influence
streamflow behavior, all of which will be applied consistently across
all models.

The target variable, daily streamflow, will be log-transformed to reduce
the influence of highly skewed flow distributions and improve model
stability. Predictions will subsequently be evaluated in both
transformed and original flow spaces to assess the implications of the
transformation and back-transformation procedure.

Prior to model fitting, basins will be split into training and test
sets. All preprocessing operations, including feature scaling and any
dimensionality reduction, will be fitted independently within each
training fold. This procedure will ensure that information from held-out
catchments does not influence model development.

The scope of this analysis will involve the usage of several machine
learning methods, both supervised and unsupervised. Ridge regression
will serve as a baseline supervised method, and it will be run on lagged
weather measurements and accuracy of nearest-gauge transfer.

After establishing a baseline ridge regression model, k-means clustering
will be applied to the static catchment attributes to identify groups of
physically similar river basins. Catchments will be classified into
groups based on their similarity across the selected standardized
attributes. The number of clusters (k) will be evaluated across three to
four candidate values. This will allow the study to examine whether the
granularity of catchment groupings influences predictive performance.
Once clustering is done, each catchment will be assigned a cluster
identifier. This identifier will subsequently serve as an alternative
representation of catchment characteristics in the supervised methods.

An important consideration is that cluster membership is derived
directly from the underlying catchment attributes. Consequently, cluster
identifiers do not introduce additional information beyond the original
attributes. Instead, they provide a potentially more parsimonious
representation of catchment similarity that may improve predictive
performance when the available training sample is small. The
hydrological coherence of the resulting clusters will also be examined
to determine whether catchments grouped together based on physical
characteristics demonstrate meaningful similarities in their observed
streamflow behavior.

Light Gradient Boosting Machine (LightGBM) will serve as the primary
supervised learning method for streamflow prediction. LightGBM is a
gradient-boosted decision tree framework that can accommodate nonlinear
relationships and interactions between meteorological and catchment
characteristics. For this model, 3 different configurations will be
compared to determine best performance: - No catchment information: The
model will use meteorological forcing variables without static catchment
attributes or cluster identifiers. - Cluster-based information: The
model will incorporate meteorological forcing variables and the
corresponding unsupervised cluster identifier derived from the k-means
clustering model. - Raw catchment attributes: The model will incorporate
meteorological forcing variables and the complete set of selected static
catchment attributes prior to log transformation.
